% !TEX root = ../../../main.tex
% Sources: end of LEC06, LEC07 photographs, and October 15, 2026 PDF.
\chapter{Eigenvalues and diagonalization}
\label{ch:linear-algebra-i:eigenvalues-and-diagonalization}

An eigenvector is a nonzero vector that a linear map sends to a scalar
multiple of itself. The characteristic polynomial detects the possible
scalars. When eigenvectors form a basis, the map has a diagonal matrix
in that basis.

Throughout this chapter, \(V\) is finite dimensional and
\(\mathbb F\) is \(\R\) or \(\C\). The choice of scalar field matters
both for the existence of eigenvalues and for diagonalization.

% !TEX root = ../../../main.tex
\section{Eigenvalues and eigenspaces}
\label{sec:linear-algebra-i:eigenvalues-and-eigenspaces}

% Sources: IMG_9267--IMG_9269, IMG_9317--IMG_9320;
% October 15 PDF, pp. 3--5.
\begin{definition}
  Let \(L:V\to V\) be linear. A scalar \(\lambda\in\mathbb F\)
  is an eigenvalue of \(L\) if there exists \(v\in V\setminus\{0\}\)
  such that
  \[
    L(v)=\lambda v.
  \]
  Such a nonzero vector \(v\) is an eigenvector for \(\lambda\).
\end{definition}

The identity map \(I\) allows us to subtract scalar multiplication from
\(L\). Since \((\lambda I)(v)=\lambda v\),
\[
  L(v)=\lambda v
  \quad\Longleftrightarrow\quad (L-\lambda I)(v)=0
  \quad\Longleftrightarrow\quad v\in\ker(L-\lambda I).
\]
For a nonzero solution, \(L-\lambda I\) is singular. Thus
\autoref{thm:linear-algebra-i:determinant-invertibility} gives
\begin{equation}
  \lambda\text{ is an eigenvalue of }L
  \quad\Longleftrightarrow\quad \det(L-\lambda I)=0.
  \label{eq:linear-algebra-i:eigenvalue-determinant-criterion}
\end{equation}
% IMG_9268 ends with a nonzero determinant. IMG_9320 and the PDF
% explicitly give zero, consistently with singularity.

\begin{definition}
  The eigenspace associated with \(\lambda\) is
  \[
    E_\lambda(L)=\ker(L-\lambda I)
      =\{v\in V\mid L(v)=\lambda v\}.
  \]
\end{definition}

An eigenspace includes the zero vector, while an eigenvector must be
nonzero. In particular, \(E_0(L)=\ker L\). A nonzero scalar multiple
of an eigenvector is another eigenvector for the same eigenvalue.

\begin{proposition}
  Every eigenspace is a subspace of \(V\). If \(\lambda\) is an
  eigenvalue, then \(\dim E_\lambda(L)>0\).
\end{proposition}

\begin{proof}
  The zero vector belongs to \(E_\lambda(L)\). For
  \(v_1,v_2\in E_\lambda(L)\) and \(a_1,a_2\in\mathbb F\),
  \begin{align*}
    L(a_1v_1+a_2v_2)
      &=a_1L(v_1)+a_2L(v_2)\\
      &=a_1\lambda v_1+a_2\lambda v_2\\
      &=\lambda(a_1v_1+a_2v_2).
  \end{align*}
  Hence \(a_1v_1+a_2v_2\in E_\lambda(L)\). For an eigenvalue,
  the eigenspace contains a nonzero vector, so its dimension is positive.
\end{proof}

\begin{proposition}
  \label{thm:linear-algebra-i:distinct-eigenspaces-intersection}
  If \(\lambda_1,\lambda_2\) are distinct eigenvalues of \(L\), then
  \[
    E_{\lambda_1}(L)\cap E_{\lambda_2}(L)=\{0\}.
  \]
\end{proposition}

\begin{proof}
  If \(v\) belongs to both eigenspaces, then
  \(L(v)=\lambda_1v\) and \(L(v)=\lambda_2v\). Subtracting gives
  \((\lambda_1-\lambda_2)v=0\). Since the scalar is nonzero,
  \(v=0\).
\end{proof}

% !TEX root = ../../hw06.tex
% Source: supplied Homework 6 PDF, page 1, Question 4.
\begin{exercise}
  \label{ex:hw06:sums-of-eigenvectors}
  Suppose, for a matrix \(A\), \(v_1\) is an eigenvector for the eigenvalue
  \(\lambda_1\) and \(v_2\) is an eigenvector for the eigenvalue
  \(\lambda_2\). What are the necessary and sufficient conditions that
  \(v_1+v_2\) be an eigenvector of \(A\) and what is its eigenvalue?
  \solutionlink{hw06:sums-of-eigenvectors}
\end{exercise}


\subsection{Matrices and coordinates}

For \(M\in\mathbb F^{n\times n}\), a nonzero column \(a\in\mathbb F^n\)
is an eigenvector with eigenvalue \(\lambda\) if \(Ma=\lambda a\).
The eigenspace is \(E_\lambda(M)=\ker(M-\lambda I)\), and the
eigenvalue criterion is \(\det(M-\lambda I)=0\).

Let \(E=(e_1,\ldots,e_n)\) be an ordered basis and let \(L(E)=EM\).
For \(v=Ea=\sum_{j=1}^n e_ja^j\),
\begin{align*}
  L(v)&=\sum_{j=1}^n a^jL(e_j)\\
    &=\sum_{j=1}^n a^j\sum_{i=1}^n e_iM^i_j\\
    &=\sum_{i=1}^n e_i\left(\sum_{j=1}^n M^i_ja^j\right)
      =E(Ma).
\end{align*}
On the other hand, \(\lambda v=E(\lambda a)\). Uniqueness of coordinates
therefore gives
\[
  L(v)=\lambda v\quad\Longleftrightarrow\quad Ma=\lambda a.
\]
Thus a nonzero vector is an eigenvector of \(L\) precisely when its
coordinate column is an eigenvector of the representing matrix.

\subsection{A motivating quotient}

% IMG_9267 also asserts that critical points are eigenvectors without a
% symmetry hypothesis. On October 22 Lucas asked to leave that claim flagged.
% Retain the supplied quotient, but do not assert the critical-point claim
% or identify its supremum with an eigenvalue until the hypothesis is clarified.
For a real matrix \(M\) and a nonzero vector \(v\in\R^n\), the quotient
\[
  \frac{v\cdot Mv}{v\cdot v}
    =\frac{\displaystyle\sum_{j,k=1}^n v^jM^j_kv^k}
      {\displaystyle\sum_{j=1}^n(v^j)^2}
\]
compares a quadratic expression with the squared length of \(v\).
The associated optimization problem is to consider
\[
  E(M)=\sup_{v\in\R^n\setminus\{0\}}
    \frac{v\cdot Mv}{v\cdot v}.
\]

% !TEX root = ../../../main.tex
\section{Characteristic polynomials and similarity}
\label{sec:linear-algebra-i:characteristic-polynomials-and-similarity}

% Sources: IMG_9269, IMG_9320--IMG_9322;
% October 15 PDF, pp. 6--10 and matrix example p. 26.
\subsection{Characteristic polynomials}

\begin{definition}
  The characteristic polynomial of \(M\in\mathbb F^{n\times n}\) is
  \[
    p_M(x)=\det(M-xI).
  \]
  For an endomorphism \(L:V\to V\), its characteristic polynomial is
  \(p_L(x)=\det(L-xI)\).
\end{definition}

Let \(\delta^i_j\) denote the entry in row \(i\), column \(j\) of
the identity matrix,
\[
  \delta^i_j=\begin{cases}1&\text{if }i=j,\\0&\text{if }i\ne j.\end{cases}
\]
The permutation formula gives
\[
  p_M(x)=\sum_{\sigma\in S_n}\varepsilon(\sigma)
    \prod_{j=1}^n\bigl(M^{\sigma(j)}_j-x\delta^{\sigma(j)}_j\bigr).
\]
This is a polynomial of degree \(n\). Its leading term is \((-x)^n\),
and its constant term is \(\det M\). Its roots in the scalar field
are exactly the eigenvalues of \(M\).

There is also a direct expansion for \(L\). Normalize a top-degree form
\(D\) on a basis \((e_1,\ldots,e_n)\). Then
\begin{align*}
  p_L(x)
    &=D((L-xI)e_1,\ldots,(L-xI)e_n)\\
    &=D(L(e_1)-xe_1,\ldots,L(e_n)-xe_n).
\end{align*}
For \(n\ge2\), multilinearity gives the expansion
\[
  \begin{aligned}
    p_L(x)&=D(L(e_1),\ldots,L(e_n))\\
      &\quad-x\sum_{j=1}^n
        D(L(e_1),\ldots,e_j,\ldots,L(e_n))\\
      &\quad+\cdots+(-x)^nD(e_1,\ldots,e_n).
  \end{aligned}
\]
In the sum displayed here, only the \(j\)-th input is replaced by
\(e_j\). The coefficient of \((-x)^{n-1}\) is
\[
  \sum_{j=1}^n D(e_1,\ldots,L(e_j),\ldots,e_n),
\]
where only the \(j\)-th input is acted on by \(L\).
Since \(D(e_1,\ldots,e_n)=1\), the leading coefficient is \((-1)^n\).
Thus
\[
  p_L(x)=(-x)^n+c_{n-1}(-x)^{n-1}+\cdots+c_1(-x)+c_0,
  \qquad c_0=\det L.
\]
If \(L(E)=EM\), then \((L-xI)(E)=E(M-xI)\), so
\autoref{thm:linear-algebra-i:determinant-representing-matrix} gives
\[
  p_L(x)=\det(L-xI)=\det(M-xI)=p_M(x).
\]

\subsection{Changing the basis}

Suppose \(E,F\) are bases with \(F=EH\), and write
\(L(E)=EM\), \(L(F)=FN\). Then
\[
  FN=L(F)=L(EH)=EMH=FH^{-1}MH,
\]
and hence \(N=H^{-1}MH\). For \(v=Ea=Fb\), the coordinate
columns satisfy \(a=Hb\) and \(b=H^{-1}a\).

\begin{definition}
  Two square matrices \(M,N\) are similar if there exists an invertible
  matrix \(H\) such that
  \[
    N=H^{-1}MH.
  \]
\end{definition}

\begin{proposition}
  \label{thm:linear-algebra-i:similarity-invariants}
  Similar matrices have the same determinant and characteristic polynomial.
  If \(N=H^{-1}MH\), then \(a\) is an eigenvector of \(M\) for
  \(\lambda\) if and only if \(H^{-1}a\) is an eigenvector of \(N\)
  for \(\lambda\).
\end{proposition}

\begin{proof}
  Multiplicativity gives
  \[
    \det N=(\det H^{-1})(\det M)(\det H)=\det M.
  \]
  Also \(N-xI=H^{-1}(M-xI)H\), so
  \[
    p_N(x)=(\det H^{-1})\det(M-xI)(\det H)=p_M(x).
  \]
  Finally,
  \[
    Ma=\lambda a
    \quad\Longleftrightarrow\quad
    N(H^{-1}a)=\lambda H^{-1}a.
  \]
  Since \(H\) is invertible, \(a\ne0\) if and only if \(H^{-1}a\ne0\).
\end{proof}

For \(\lambda=0\), this gives an isomorphism of kernels. Its two directions
are shown in \autoref{fig:linear-algebra-i:similar-matrix-kernels},
\[
  \ker M\longrightarrow\ker N,\quad a\longmapsto H^{-1}a,
  \qquad b\longmapsto Hb\text{ in the reverse direction}.
\]
% IMG_9322 labels the direction of its H^{-1} arrow inconsistently.
% The calculation and N=H^{-1}MH fix the directions used here.
\begin{marginfigure}
  \centering
  % Editable version of the kernel arrows in IMG_9322.
% Convention: N=H^{-1}MH. The inverse directions follow the displayed algebra.
\begin{tikzpicture}[
  x=1cm,y=1cm,
  every node/.style={font=\footnotesize},
  >={Latex[length=1.6mm,width=1.1mm]},
  line width=0.45pt
]
  \node (m) at (0,0) {\(\ker M\)};
  \node (n) at (2.6,0) {\(\ker N\)};
  \draw[->] (m.north east) to[bend left=24]
    node[above=3pt] {\(a\mapsto H^{-1}a\)} (n.north west);
  \draw[->] (n.south west) to[bend left=24]
    node[below=3pt] {\(b\mapsto Hb\)} (m.south east);
\end{tikzpicture}

  \caption[Kernel isomorphism under similarity]{For \(N=H^{-1}MH\),
    the changes of coordinates restrict to inverse isomorphisms of kernels.}
  \label{fig:linear-algebra-i:similar-matrix-kernels}
\end{marginfigure}

\begin{example}
  \label{ex:linear-algebra-i:same-polynomial-not-similar}
  The converse to characteristic-polynomial invariance fails. Let
  \[
    A=\begin{bmatrix}0&0\\0&0\end{bmatrix},\qquad
    N=\begin{bmatrix}0&1\\0&0\end{bmatrix}.
  \]
  Both characteristic polynomials are \(x^2\), but
  \[
    \ker A=\mathbb F^2,\qquad
    \ker N=\left\{\begin{bmatrix}t\\0\end{bmatrix}\;\middle|\;t\in\mathbb F\right\}.
  \]
  Their kernels have different dimensions, so they cannot be similar.
\end{example}

% !TEX root = ../../../main.tex
\section{Computing eigenvalues and eigenspaces}
\label{sec:linear-algebra-i:computing-eigenspaces}

% Sources: IMG_9323--IMG_9326; October 15 PDF, pp. 11--16.
\subsection{Diagonal and triangular matrices}

\begin{example}
  For
  \[
    D=\begin{bmatrix}a&0&0\\0&b&0\\0&0&c\end{bmatrix},
  \]
  the standard basis vectors satisfy
  \(De_1=ae_1\), \(De_2=be_2\), and \(De_3=ce_3\).
  The characteristic polynomial is
  \[
    p_D(x)=\det\begin{bmatrix}
      a-x&0&0\\0&b-x&0\\0&0&c-x
    \end{bmatrix}=(a-x)(b-x)(c-x).
  \]
  Thus the eigenvalues are the diagonal entries.
\end{example}

A triangular matrix also has its diagonal entries as eigenvalues, since
subtracting \(xI\) preserves triangularity. For instance,
\[
  T=\begin{bmatrix}a&0&0\\b&c&0\\d&e&f\end{bmatrix}
  \quad\Longrightarrow\quad
  p_T(x)=(a-x)(c-x)(f-x).
\]
Here \(Te_3=fe_3\), so \(e_3\) is an eigenvector for \(f\).
Finding eigenvectors for \(a\) and \(c\) requires solving the corresponding
kernel equations. In general, for triangular \(T=(T^i_j)\),
\[
  \det T=T^1_1\cdots T^n_n,\qquad
  p_T(x)=\prod_{j=1}^n(T^j_j-x).
\]

% !TEX root = ../../hw06.tex
% Source: supplied Homework 6 PDF, page 1, Question 1.
\begin{exercise}
  \label{ex:hw06:triangular-matrix-eigenvectors}
  Find all eigenvectors of the matrix
  \[
    M=\begin{bmatrix}
      1&1&0\\
      0&2&1\\
      0&0&2
    \end{bmatrix}.
  \]
  \solutionlink{hw06:triangular-matrix-eigenvectors}
\end{exercise}


% !TEX root = ../../hw06.tex
% Source: supplied Homework 6 PDF, page 1, Question 2.
% The source calls L a map and implicitly uses its linearity; do not add a
% hypothesis or a solution without the author's clarification.
\begin{exercise}
  \label{ex:hw06:eigenvectors-in-a-basis}
  Let \((e_1,e_2,e_3)\) be a basis of \(V\) and \(L:V\to V\) be the map
  given by
  \[
    \begin{aligned}
      L(e_1)&=e_1+e_2\\
      L(e_2)&=2e_2+e_3\\
      L(e_3)&=2e_3.
    \end{aligned}
  \]
  Find the eigenvectors of \(L\) by solving
  \[
    L(v)=\lambda v,
  \]
  where
  \[
    v=e_1a^1+e_2a^2+e_3a^3,
  \]
  directly for \(a^1,a^2,a^3,\lambda\).
  \solutionlink{hw06:eigenvectors-in-a-basis}
\end{exercise}


\subsection{A matrix that interchanges coordinates}

\begin{example}
  \label{ex:linear-algebra-i:coordinate-swap-eigenspaces}
  Let
  \[
    A=\begin{bmatrix}0&1\\1&0\end{bmatrix}.
  \]
  Its characteristic polynomial is
  \(p_A(x)=x^2-1=(x-1)(x+1)\). Row reduction gives
  \[
    A-\lambda I=
    \begin{bmatrix}-\lambda&1\\1&-\lambda\end{bmatrix}
    \sim\begin{bmatrix}1&-\lambda\\-\lambda&1\end{bmatrix}
    \sim\begin{bmatrix}1&-\lambda\\0&1-\lambda^2\end{bmatrix}.
  \]
  If \(1-\lambda^2\ne0\), the matrix reduces to the identity and its
  kernel is zero. If \(\lambda=1\), the kernel equation is
  \(v^1-v^2=0\), giving
  \[
    E_1(A)=\left\{\begin{bmatrix}t\\t\end{bmatrix}\;\middle|\;t\in\mathbb F\right\}
      =\operatorname{span}\left\{\begin{bmatrix}1\\1\end{bmatrix}\right\}.
  \]
  If \(\lambda=-1\), the equation is \(v^1+v^2=0\), giving
  \[
    E_{-1}(A)=\left\{\begin{bmatrix}t\\-t\end{bmatrix}\;\middle|\;t\in\mathbb F\right\}
      =\operatorname{span}\left\{\begin{bmatrix}1\\-1\end{bmatrix}\right\}.
  \]
  The nonzero vectors in these spaces are the eigenvectors.
  Over \(\R\), their two directions appear in
  \autoref{fig:linear-algebra-i:swap-eigendirections}.
  \begin{marginfigure}
    \centering
    % Exact eigendirections of A=[0 1;1 0] over R.
\begin{tikzpicture}[
  x=10mm,y=10mm,font=\footnotesize,line width=0.45pt,
  vector/.style={-{Latex[length=1.6mm,width=1.1mm]}},
  axis/.style={vector,black!30,line width=0.35pt},
  direction/.style={dashed,black!55,line width=0.35pt}
]
  \draw[axis] (-1.7,0) -- (1.8,0) node[right=2pt] {\(x\)};
  \draw[axis] (0,-1.6) -- (0,1.8) node[above=2pt] {\(y\)};
  \draw[direction] (-1.5,-1.5) -- (1.5,1.5) node[above right=1pt] {\(E_1\)};
  \draw[direction] (-1.5,1.5) -- (1.5,-1.5) node[below right=1pt] {\(E_{-1}\)};
  \draw[vector] (0,0) -- (1,1) node[right=3pt] {\(v=Av\)};
  \draw[vector] (0,0) -- (1,-1) node[right=3pt] {\(w\)};
  \draw[vector] (0,0) -- (-1,1) node[left=3pt] {\(Aw\)};
  \node[below left=2pt] at (0,0) {\(0\)};
  \path (current bounding box.south) ++(0,-2mm);
\end{tikzpicture}

    \caption[Eigendirections of the coordinate swap]{The map
      \(A(x,y)=(y,x)\) fixes the line \(E_1\) and reverses vectors on
      \(E_{-1}\). Here \(v=(1,1)\), \(w=(1,-1)\), and \(Aw=-w\).}
    \label{fig:linear-algebra-i:swap-eigendirections}
  \end{marginfigure}
\end{example}
% IMG_9325's final vector for -1 has a sign inconsistent with v1=-v2.
% Use the preceding kernel equation and PDF p.15.

\subsection{A matrix with no real eigenvalue}

\begin{example}
  \label{ex:linear-algebra-i:rotation-no-real-eigenvalue}
  Consider
  \[
    B=\begin{bmatrix}0&-1\\1&0\end{bmatrix},\qquad
    B\begin{bmatrix}x\\y\end{bmatrix}=\begin{bmatrix}-y\\x\end{bmatrix}.
  \]
  Its characteristic polynomial is \(p_B(x)=x^2+1\), and
  \[
    B-\lambda I=
    \begin{bmatrix}-\lambda&-1\\1&-\lambda\end{bmatrix}
    \sim\begin{bmatrix}1&-\lambda\\-\lambda&-1\end{bmatrix}
    \sim\begin{bmatrix}1&-\lambda\\0&-1-\lambda^2\end{bmatrix}.
  \]
  For real \(\lambda\), the last diagonal entry is strictly negative,
  so the kernel is zero. Thus \(B\) has no real eigenvalues.
  Over \(\C\), the roots of \(x^2+1\) are \(i\) and \(-i\).
  Geometrically, \(B\) makes a counterclockwise quarter-turn in the
  real plane, as in
  \autoref{fig:linear-algebra-i:rotation-no-real-eigendirection}.
  \begin{marginfigure}
    \centering
    % Exact action B(x,y)=(-y,x). The sample direction is schematic, not a new example.
\begin{tikzpicture}[
  x=9mm,y=9mm,font=\footnotesize,line width=0.45pt,
  vector/.style={-{Latex[length=1.6mm,width=1.1mm]}},
  axis/.style={vector,black!35,line width=0.35pt}
]
  \coordinate (O) at (0,0);
  \coordinate (V) at (1.8,0.6);
  \coordinate (BV) at (-0.6,1.8);
  \draw[axis] (-1.1,0) -- (2.4,0) node[right=2pt] {\(x\)};
  \draw[axis] (0,-0.5) -- (0,2.4) node[above=2pt] {\(y\)};
  \draw[vector] (O) -- (V) node[right=3pt] {\(v\)};
  \draw[vector] (O) -- (BV) node[above left=3pt] {\(Bv\)};
  \draw[vector,black!60] (18.43495:1.1) arc[start angle=18.43495,end angle=108.43495,radius=1.1];
  \node at (63.43495:0.7) {\(90^\circ\)};
  \node[below left=2pt] at (O) {\(0\)};
  \path (current bounding box.south) ++(0,-2mm);
\end{tikzpicture}

    \caption[A quarter-turn in the real plane]{For every real nonzero
      \(v\), the vector \(Bv\) is perpendicular to \(v\), so it cannot
      be a real scalar multiple of \(v\).}
    \label{fig:linear-algebra-i:rotation-no-real-eigendirection}
  \end{marginfigure}
\end{example}

% !TEX root = ../../../main.tex
\section{Complex eigenvalues}
\label{sec:linear-algebra-i:complex-eigenvalues}

% Sources: IMG_9326--IMG_9327; October 15 PDF, pp. 17--21.
\subsection{Complex arithmetic}

A complex number has the form \(z=x+iy\), with \(x,y\in\R\) and
\(i^2=-1\). Its real and imaginary parts are \(x\) and \(y\), respectively.
\begin{marginfigure}
  \centering
  % Generic complex number and its conjugate; the positive quadrant is illustrative.
\begin{tikzpicture}[
  x=10mm,y=10mm,font=\footnotesize,line width=0.45pt,
  axis/.style={black!40,line width=0.35pt,-{Latex[length=1.6mm,width=1.1mm]}},
  guide/.style={dashed,black!45,line width=0.35pt}
]
  \coordinate (Z) at (1.7,1.3);
  \coordinate (C) at (1.7,-1.3);
  \draw[axis] (-0.4,0) -- (2.8,0) node[right=2pt] {\(\operatorname{Re}\)};
  \draw[axis] (0,-1.8) -- (0,1.9) node[above=2pt] {\(\operatorname{Im}\)};
  \draw[guide] (0,1.3) -- (Z) -- (C) -- (0,-1.3);
  \draw (0,0) -- (Z) node[midway,above left=1pt] {\(\lvert z\rvert\)};
  \draw[black!55] (0,0) -- (C);
  \fill (Z) circle[radius=1.5pt];
  \fill (C) circle[radius=1.5pt];
  \node[above right=2pt] at (Z) {\(z=x+iy\)};
  \node[below right=2pt] at (C) {\(\overline z=x-iy\)};
  \node[left=3pt] at (0,1.3) {\(y\)};
  \node[left=3pt] at (0,-1.3) {\(-y\)};
  \node[above right=1pt] at (1.7,0) {\(x\)};
  \node[below left=2pt] at (0,0) {\(0\)};
  \path (current bounding box.south) ++(0,-2mm);
\end{tikzpicture}

  \caption[Conjugation in the complex plane]{A schematic for \(x,y>0\).
    Conjugation reflects a point across the real axis. The distance from
    zero is \(\lvert z\rvert=\sqrt{x^2+y^2}\) for both points.}
  \label{fig:linear-algebra-i:complex-conjugation}
\end{marginfigure}
Powers of \(i\) repeat with period four,
\[
  i^2=-1,\qquad i^3=-i,\qquad i^4=1,\qquad i^5=i,\qquad i^6=-1.
\]
The conjugate of \(z\) is \(\overline z=x-iy\), and its modulus is
the nonnegative number satisfying
\[
  |z|^2=z\overline z=x^2+y^2.
\]
For \(z\ne0\),
\[
  \frac1z=\frac{\overline z}{|z|^2}=\frac{x-iy}{x^2+y^2}.
\]
More generally, division by \(w\ne0\) can be written as
\(z/w=z\overline w/|w|^2\).

\autoref{fig:linear-algebra-i:complex-conjugation} identifies \(z\)
with the point \((x,y)\) and its conjugate with \((x,-y)\).

\subsection{Returning to the rotation example}

Regard the real matrix \(B\) from
\autoref{ex:linear-algebra-i:rotation-no-real-eigenvalue} as a matrix
acting on \(\C^2\). For \(\lambda=i\), its kernel equations give
\[
  v^1-iv^2=0\quad\Longleftrightarrow\quad v^2=-iv^1.
\]
For \(\lambda=-i\), they give \(v^1+iv^2=0\). Consequently,
\[
  \begin{aligned}
    E_i(B)&=\left\{\begin{bmatrix}t\\-it\end{bmatrix}\;\middle|\;t\in\C\right\},\\
    E_{-i}(B)&=\left\{\begin{bmatrix}t\\it\end{bmatrix}\;\middle|\;t\in\C\right\}.
  \end{aligned}
\]
In each case, the eigenvectors are obtained by taking \(t\ne0\).
% IMG_9327 has an isolated +1 in B-iI. The original matrix, its stated
% action, both kernel equations, and PDF p.18 fix that entry as -1.

% !TEX root = ../../hw06.tex
% Source: supplied Homework 6 PDF, page 1, Question 5.
\begin{exercise}
  \label{ex:hw06:rotation-eigenvalues}
  For any \(\alpha\in\R\), the real matrix
  \[
    R_\alpha=\begin{bmatrix}
      \cos\alpha&-\sin\alpha\\
      \sin\alpha&\cos\alpha
    \end{bmatrix}
  \]
  is called a rotation matrix, because it rotates each vector counterclockwise
  by the angle \(\alpha\). Compute the (possibly complex) eigenvalues and
  eigenvectors of \(R_\alpha\)
  \solutionlink{hw06:rotation-eigenvalues}
\end{exercise}


\subsection{Roots and multiplicities}

\begin{theorem}[Fundamental theorem of algebra]
  Every complex polynomial of degree \(n\ge1\),
  \[
    p(z)=a_nz^n+\cdots+a_0,\qquad a_n\ne0,
  \]
  factors as
  \[
    p(z)=a_n(z-\lambda_1)\cdots(z-\lambda_n),
  \]
  where \(\lambda_1,\ldots,\lambda_n\in\C\), with repetitions allowed.
\end{theorem}
% PDF p.19 states this theorem without a proof.

The characteristic polynomial of an \(n\)-by-\(n\) complex matrix
therefore has \(n\) roots counted with multiplicity. The multiplicity
of \(\lambda\) is the number of factors equal to \((z-\lambda)\).
The \(n\) eigenvalues need not be distinct. Since the convention here is
\(p_M(z)=\det(M-zI)\), its factorization has leading coefficient
\((-1)^n\).

\begin{example}
  For the identity matrix on \(\mathbb F^2\),
  \[
    p_I(x)=(1-x)^2,\qquad E_1(I)=\mathbb F^2.
  \]
  Its only eigenvalue is \(1\), with multiplicity two, and every
  nonzero vector is an eigenvector.

  In contrast, let
  \[
    J=\begin{bmatrix}1&1\\0&1\end{bmatrix}.
  \]
  It has the same characteristic polynomial, but
  \[
    (J-I)\begin{bmatrix}v^1\\v^2\end{bmatrix}
      =\begin{bmatrix}v^2\\0\end{bmatrix}.
  \]
  Hence
  \[
    E_1(J)=\left\{\begin{bmatrix}t\\0\end{bmatrix}\;\middle|\;t\in\mathbb F\right\},
  \]
  which has dimension one.
\end{example}
% PDF p.21 prints (0,c) after deriving v2=0. Use (c,0), as its
% displayed multiplication and the definition of the kernel require.

The different actions of \(I\) and \(J\) over \(\R\) are pictured in
\autoref{fig:linear-algebra-i:identity-and-shear-eigenspaces}.
\begin{figure}[htbp]
  \centering
  % The supplied matrices I and J=[1 1;0 1]. Equal units in both panels.
\begin{tikzpicture}[
  x=13mm,y=13mm,font=\footnotesize,line width=0.45pt,
  vector/.style={-{Latex[length=1.6mm,width=1.1mm]}},
  axis/.style={vector,black!40,line width=0.35pt},
  guide/.style={dashed,black!45,line width=0.35pt}
]
  \begin{scope}
    \node at (1.1,2) {\(I\)};
    \fill[black!8] (0,0) -- (1,0) -- (1,1) -- (0,1) -- cycle;
    \draw[axis] (-0.4,0) -- (2.6,0) node[right=2pt] {\(x\)};
    \draw[axis] (0,-0.3) -- (0,1.6) node[above=2pt] {\(y\)};
    \draw (1,0) -- (1,1) -- (0,1);
    \draw[vector] (0,0) -- (1,0) node[below=3pt] {\(e_1\)};
    \draw[vector] (0,0) -- (0,1) node[left=3pt] {\(e_2\)};
    \node at (1.1,-0.9) {\(E_1(I)=\R^2\)};
  \end{scope}
  \begin{scope}[xshift=61mm]
    \node at (1.1,2) {\(J\)};
    \fill[black!8] (0,0) -- (1,0) -- (2,1) -- (1,1) -- cycle;
    \draw[axis] (-0.4,0) -- (2.6,0) node[right=2pt] {\(x\)};
    \draw[axis] (0,-0.3) -- (0,1.6) node[above=2pt] {\(y\)};
    \draw[guide] (0,0) -- (0,1) -- (1,1) -- (1,0);
    \draw (1,0) -- (2,1) -- (1,1);
    \draw[vector] (0,0) -- (1,0) node[below=3pt] {\(Je_1=e_1\)};
    \draw[vector] (0,0) -- (1,1) node[above left=3pt] {\(Je_2\)};
    \node at (1.1,-0.9) {\(E_1(J)=\operatorname{span}\{e_1\}\)};
  \end{scope}
\end{tikzpicture}

  \caption[Equal multiplicities, different eigenspaces]{The identity
    fixes every vector. The shear \(J\) fixes only the horizontal axis,
    although both characteristic polynomials are \((1-x)^2\).
    The dashed unit square on the right shows the shape before applying
    \(J\); the shaded parallelogram is its image.}
  \label{fig:linear-algebra-i:identity-and-shear-eigenspaces}
\end{figure}

\begin{proposition}
  If \(\lambda\) is an eigenvalue of a complex matrix and has
  multiplicity \(m_\lambda\) in its characteristic polynomial, then
  \[
    1\le\dim E_\lambda\le m_\lambda.
  \]
\end{proposition}
% PDF p.21 states the upper bound without supplying a proof.

% !TEX root = ../../hw06.tex
% Source: supplied Homework 6 PDF, page 1, Question 3, parts 3.1--3.4.
% Preserve the unrestricted field in 3.2 and the absent field in 3.1;
% the truth values are for the reader to assess, not assumptions to repair.
\begin{exercise}
  \label{ex:hw06:eigenvalue-true-or-false}
  True or false? Explain. If false, provide a counterexample.
  \begin{enumerate}[{3}.1.][4]
    \item If a matrix has one eigenvector, then it has infinitely many
      eigenvectors.
    \item Every matrix over a field \(\mathbb F\) has at least one eigenvalue
      in \(\mathbb F\).
    \item Similar matrices always have the same eigenvectors.
    \item If two matrices have the same eigenvalues, then they are similar
  \end{enumerate}
  \solutionlink{hw06:eigenvalue-true-or-false}
\end{exercise}


% !TEX root = ../../../main.tex
\section{Diagonalization}
\label{sec:linear-algebra-i:diagonalization}

% Sources: IMG_9268, IMG_9323--IMG_9324, IMG_9328;
% October 15 PDF, pp. 22--26.
\begin{definition}
  A square matrix \(M\) is diagonalizable over \(\mathbb F\) if it is
  similar to a diagonal matrix. Equivalently, there exist an invertible
  matrix \(S\) and a diagonal matrix \(D\) such that
  \[
    M=SDS^{-1}.
  \]
\end{definition}

\begin{theorem}[Diagonalization and eigenbases]
  \label{thm:linear-algebra-i:diagonalization-eigenbasis}
  A matrix \(M\in\mathbb F^{n\times n}\) is diagonalizable if and only
  if \(\mathbb F^n\) has a basis of eigenvectors of \(M\).
\end{theorem}

\begin{marginfigure}
  \centering
  % MS=SD: S converts eigenbasis coordinates into the original coordinates.
\begin{tikzpicture}[
  x=1mm,y=1mm,font=\footnotesize,line width=0.45pt,
  every node/.style={inner sep=3pt},
  map/.style={-{Latex[length=1.6mm,width=1.1mm]}}
]
  \node (c) at (0,24) {\(c\)};
  \node (Dc) at (32,24) {\(Dc\)};
  \node (Sc) at (0,0) {\(Sc\)};
  \node (SDc) at (32,0) {\(SDc\)};
  \draw[map] (c) -- node[above] {\(D\)} (Dc);
  \draw[map] (c) -- node[left] {\(S\)} (Sc);
  \draw[map] (Dc) -- node[right] {\(S\)} (SDc);
  \draw[map] (Sc) -- node[below] {\(M\)} (SDc);
  \path (current bounding box.south) ++(0,-2mm);
\end{tikzpicture}

  \caption[Diagonalization through coordinates]{The top row acts in
    eigenbasis coordinates; the bottom acts in the original coordinates.
    Both routes give \(MSc=SDc\).}
  \label{fig:linear-algebra-i:diagonalization-coordinates}
\end{marginfigure}

\begin{proof}
  Suppose \(M=SDS^{-1}\), where \(D\) has diagonal entries
  \(\lambda_1,\ldots,\lambda_n\). Let \(c_k=Se_k\) be column \(k\)
  of \(S\). Since \(S\) is invertible, its columns form a basis, and
  \begin{align*}
    Mc_k&=MSe_k=SDS^{-1}Se_k\\
      &=SDe_k=S(\lambda_ke_k)=\lambda_kc_k.
  \end{align*}
  Thus each column is an eigenvector.

  Conversely, suppose \((b_1,\ldots,b_n)\) is a basis of eigenvectors,
  with \(Mb_k=\lambda_kb_k\). Let \(S\) have these basis vectors as
  columns, and let \(D\) have the corresponding eigenvalues on its diagonal.
  The column equations say \(MS=SD\). Since \(S\) is invertible,
  \(M=SDS^{-1}\).
\end{proof}

The identity \(MS=SD\) is visualized in
\autoref{fig:linear-algebra-i:diagonalization-coordinates}.
If \(c\) is an eigenbasis coordinate column, \(S\) reconstructs the vector
\(Sc\), while \(D\) multiplies each coordinate by its eigenvalue.

\begin{definition}
  An endomorphism \(L:V\to V\) is diagonalizable if \(V\) has a basis
  \((e_1,\ldots,e_n)\) of eigenvectors of \(L\).
\end{definition}

For such a basis,
\[
  L(e_k)=\lambda_ke_k,\qquad
  L(E)=E\begin{bmatrix}
    \lambda_1&0&\cdots&0\\
    0&\lambda_2&\cdots&0\\
    \vdots&\vdots&\ddots&\vdots\\
    0&0&\cdots&\lambda_n
  \end{bmatrix}.
\]
Equivalently, the matrix representing \(L\) in any basis is similar
to a diagonal matrix.

\subsection{Failure of diagonalization}

\begin{example}
  Let \((e_1,e_2)\) be a basis of \(V\), and define
  \[
    L(e_1)=e_2,\qquad L(e_2)=0.
  \]
  For \(v=ae_1+be_2\), \(L(v)=ae_2\). The eigenvector equation is
  \[
    ae_2=\lambda(ae_1+be_2).
  \]
  If \(\lambda\ne0\), comparison of coefficients gives \(a=0\)
  and then \(b=0\), so there is no nonzero eigenvector. For
  \(\lambda=0\), it gives \(a=0\). Thus the only eigenvalue is zero
  and its eigenspace is \(\operatorname{span}\{e_2\}\). It cannot
  provide a basis of \(V\), so \(L\) is not diagonalizable.
\end{example}

The nonzero matrix \(N\) in
\autoref{ex:linear-algebra-i:same-polynomial-not-similar} similarly has
only the eigenvalue zero and a one-dimensional eigenspace, so it is not
diagonalizable. The rotation matrix in
\autoref{ex:linear-algebra-i:rotation-no-real-eigenvalue} is not
diagonalizable over \(\R\), since it has no real eigenvectors.

% !TEX root = ../../hw06.tex
% Source: supplied Homework 6 PDF, page 2, Question 6.
% Retain the source's "zero matrix" wording for the linear transformation;
% its unstated nonzero-space hypothesis is recorded in the README.
\begin{exercise}
  \label{ex:hw06:nilpotent-eigenvalues}
  A linear transformation \(A:V\to V\) is called \emph{nilpotent} if for some
  positive integer \(k\), \(A^k\) is the zero matrix. Prove that the only
  eigenvalue of \(A\) is \(0\).
  \solutionlink{hw06:nilpotent-eigenvalues}
\end{exercise}


% !TEX root = ../../../main.tex
\section{Direct sums and diagonalization criteria}
\label{sec:linear-algebra-i:direct-sums}

% Sources: IMG_9328--IMG_9330; October 22 handwritten PDF, pp. 1--3.
% The October 15 PDF does not contain this material.
\subsection{Transversal subspaces}

Suppose \(S,T\) are subspaces of \(V\) and \(S\cap T=\{0\}\).
If \(s\in S\), \(t\in T\), and \(s-t=0\), then
\(s=t\in S\cap T\), so both vectors are zero.

% !TEX root = ../../hw07.tex
% Source: ma-gy7033fall2026homework7 (1).pdf, p. 1, Problem 3.
% "Without relying on the lectures, prove" becomes "Prove directly"
% so that the shared statement is self-contained in the book.
\begin{exercise}
  \label{ex:hw07:disjoint-subspaces-independence}
  (10 points) Prove directly that if \(V_1,V_2\) are subspaces of \(V\) and
  \[
    V_1\cap V_2=\{0\},
  \]
  then if \(v_1\in V_1\) and \(v_2\in V_2\) are nonzero vectors,
  then \(\{v_1,v_2\}\) is a linearly independent set.
  \solutionlink{hw07:disjoint-subspaces-independence}
\end{exercise}


For a family \(S_1,\ldots,S_k\), write
\[
  S_1+\cdots+S_k
    =\{s_1+\cdots+s_k\mid s_j\in S_j\}.
\]
This is the span of the subspaces together.

\begin{definition}
  A family of subspaces \(S_1,\ldots,S_k\) is transversal if
  \[
    S_j\cap\sum_{i\ne j}S_i=\{0\}\qquad\text{for every }j.
  \]
\end{definition}
% Retain the course's defined use of transversal. Do not replace the
% family condition by pairwise trivial intersections, which is weaker.

The coordinate axes are transversal. In contrast, three distinct lines
through zero in a plane are not transversal, as illustrated in
\autoref{fig:linear-algebra-i:three-lines-not-transversal}.

\begin{marginfigure}[3\baselineskip]
  \centering
  % The October 22 p. 1 example: three distinct lines in the real plane.
\begin{tikzpicture}[
  x=10mm,y=10mm,font=\footnotesize,line width=0.45pt,
  continued/.style={{Latex[length=1.6mm,width=1.1mm]}-{Latex[length=1.6mm,width=1.1mm]}}
]
  \draw[continued,black!55] (-1.7,0) -- (1.7,0) node[right=2pt] {\(S_1\)};
  \draw[continued,black!55] (0,-1.5) -- (0,1.7) node[above=2pt] {\(S_2\)};
  \draw[continued] (-1.4,-1.05) -- (1.5,1.125) node[above right=2pt] {\(S_3\)};
  \fill (0,0) circle[radius=1.3pt];
  \node[below right=3pt] at (0,0) {\(0\)};
  \path (current bounding box.south) ++(0,-2mm);
\end{tikzpicture}

  \caption[Three lines and transversality]{A real-plane schematic.
    Each pair of distinct lines meets only at zero, but any two span the
    plane. Thus \(S_3\cap(S_1+S_2)=S_3\ne\{0\}\).}
  \label{fig:linear-algebra-i:three-lines-not-transversal}
\end{marginfigure}

\begin{proposition}
  \label{thm:linear-algebra-i:transversal-unique-components}
  Let \(S_1,\ldots,S_k\) be transversal, with \(s_j\in S_j\).
  If \(s_1+\cdots+s_k=0\), then every \(s_j=0\).
  Consequently, if also \(t_j\in S_j\) and
  \(s_1+\cdots+s_k=t_1+\cdots+t_k\), then \(s_j=t_j\) for every \(j\).
\end{proposition}

\begin{proof}
  Isolate the last component in the zero relation. It satisfies
  \[
    s_k=-\sum_{j=1}^{k-1}s_j
      \in S_k\cap(S_1+\cdots+S_{k-1})=\{0\}.
  \]
  Hence \(s_k=0\). Repeating with the remaining components gives
  \(s_1=\cdots=s_k=0\).

  For two equal sums, subtract them. Each \(s_j-t_j\) belongs to
  \(S_j\), and their sum is zero. The first conclusion gives
  \(s_j-t_j=0\) for every \(j\).
\end{proof}

\begin{definition}
  If \(S_1,\ldots,S_k\) are transversal and
  \(S_1+\cdots+S_k=V\), then \(V\) is their direct sum, written
  \[
    V=S_1\oplus\cdots\oplus S_k=\bigoplus_{j=1}^k S_j.
  \]
\end{definition}

\subsection{Combining bases}

For two complementary lines in \(\R^2\), each vector has one component
on each line, as in \autoref{fig:linear-algebra-i:direct-sum-components}.
\begin{marginfigure}
  \centering
  % Two-dimensional real schematic; the two subspace lines span the plane.
\begin{tikzpicture}[
  x=11mm,y=11mm,font=\footnotesize,line width=0.45pt,
  vector/.style={-{Latex[length=1.6mm,width=1.1mm]}},
  continued/.style={black!50,line width=0.35pt,{Latex[length=1.6mm,width=1.1mm]}-{Latex[length=1.6mm,width=1.1mm]}},
  guide/.style={dashed,black!45,line width=0.35pt}
]
  \coordinate (O) at (0,0);
  \coordinate (S) at (1.8,0.3);
  \coordinate (T) at (0.35,1.4);
  \coordinate (V) at ($(S)+(T)$);
  \draw[continued] (-0.6,-0.1) -- (2.7,0.45) node[right=2pt] {\(S\)};
  \draw[continued] (-0.15,-0.6) -- (0.6,2.4) node[above=2pt] {\(T\)};
  \draw[guide] (S) -- (V) -- (T);
  \draw[vector] (O) -- (S) node[midway,below=3pt] {\(s\)};
  \draw[vector] (O) -- (T) node[midway,left=3pt] {\(t\)};
  \draw[vector] (O) -- (V) node[above right=2pt] {\(v\)};
  \node[below left=2pt] at (O) {\(0\)};
  \path (current bounding box.south) ++(0,-2mm);
\end{tikzpicture}

  \caption[Components in a direct sum]{A real-plane schematic of
    \(V=S\oplus T\). The decomposition \(v=s+t\), with \(s\in S\)
    and \(t\in T\), is unique. The summands need not be perpendicular.}
  \label{fig:linear-algebra-i:direct-sum-components}
\end{marginfigure}

\begin{proposition}
  \label{thm:linear-algebra-i:direct-sum-bases}
  If \(V=S_1\oplus\cdots\oplus S_k\), concatenating a basis of each
  \(S_j\) gives a basis of \(V\). In particular,
  \[
    \dim V=\sum_{j=1}^k\dim S_j.
  \]
\end{proposition}

\begin{proof}
  Let \((e_{j,1},\ldots,e_{j,n_j})\) be a basis of \(S_j\).
  Suppose a linear combination of all these vectors is zero,
  \[
    \sum_{j=1}^k\sum_{\ell=1}^{n_j}a_{j,\ell}e_{j,\ell}=0.
  \]
  Group the terms into vectors
  \(s_j=\sum_{\ell=1}^{n_j}a_{j,\ell}e_{j,\ell}\in S_j\).
  Transversality gives \(s_1=\cdots=s_k=0\). Since each group is a
  linear combination of a basis of \(S_j\), all its coefficients
  \(a_{j,\ell}\) are zero. Thus the concatenated list is independent.

  The subspaces together span \(V\), and the basis vectors in each
  group span its subspace. Hence the concatenated list spans \(V\).
  It is therefore a basis, with \(n_1+\cdots+n_k\) vectors.
\end{proof}

\subsection{Transversality of distinct eigenspaces}

\begin{theorem}[Transversal eigenspaces]
  \label{thm:linear-algebra-i:distinct-eigenspaces-transversal}
  Let \(\lambda_1,\ldots,\lambda_k\) be distinct eigenvalues of
  \(L:V\to V\). Their eigenspaces
  \(E_{\lambda_1}(L),\ldots,E_{\lambda_k}(L)\) are transversal.
\end{theorem}

\begin{proof}
  We show that a relation
  \[
    v_1+\cdots+v_k=0,\qquad v_j\in E_{\lambda_j}(L),
  \]
  forces every \(v_j=0\). The argument proceeds by induction on the
  number of eigenspaces. For one eigenspace, the relation already says
  \(v_1=0\).

  Suppose the result holds for \(k-1\) distinct eigenspaces. Apply
  \(L\) to the relation and subtract \(\lambda_k\) times the original
  relation. Since \(L(v_j)=\lambda_jv_j\), this gives
  \[
    \sum_{j=1}^{k-1}(\lambda_j-\lambda_k)v_j=0.
  \]
  The \(k\)-th term has disappeared. Each remaining vector
  \((\lambda_j-\lambda_k)v_j\) still belongs to \(E_{\lambda_j}(L)\).
  By the induction hypothesis, all these vectors are zero. The scalars
  \(\lambda_j-\lambda_k\) are nonzero, so \(v_j=0\) for \(j<k\).
  The original relation now gives \(v_k=0\).

  A vector in one eigenspace that is a sum of vectors in the others
  produces such a zero relation by moving all terms to one side.
  The result just proved forces that vector to be zero. This is
  transversality.
\end{proof}
% October 22 p. 2 supplies this elimination argument and says to repeat it.
% The induction spells out that supplied iteration, replacing p. 1's
% explicitly abandoned attempt. Pairwise intersections alone are insufficient.

\subsection{Criteria for an eigenbasis}

\begin{theorem}[Spanning eigenspaces and diagonalization]
  \label{thm:linear-algebra-i:spanning-eigenspaces-diagonalizable}
  If \(\lambda_1,\ldots,\lambda_k\) are eigenvalues of \(L:V\to V\)
  and
  \[
    E_{\lambda_1}(L)+\cdots+E_{\lambda_k}(L)=V,
  \]
  then \(L\) is diagonalizable.
\end{theorem}

\begin{proof}
  Keep one copy of each distinct eigenspace in the sum. By
  \autoref{thm:linear-algebra-i:distinct-eigenspaces-transversal}, these
  eigenspaces are transversal, so the spanning hypothesis gives a
  direct-sum decomposition of \(V\). Choose a basis in each eigenspace.
  By \autoref{thm:linear-algebra-i:direct-sum-bases}, their concatenation
  is a basis of \(V\). Every vector in this basis is an eigenvector.
  Thus \(L\) is diagonalizable.
\end{proof}

The converse also holds. An endomorphism is diagonalizable exactly when
its eigenspaces span \(V\).
% October 22 p. 3 asserts the converse without writing a separate proof.

\begin{corollary}
  \label{thm:linear-algebra-i:distinct-eigenvalues-diagonalizable}
  If \(\dim V=n\) and \(L:V\to V\) has \(n\) distinct eigenvalues
  in \(\mathbb F\), then \(L\) is diagonalizable over \(\mathbb F\).
  In particular, a two-by-two matrix with two distinct eigenvalues in
  \(\mathbb F\) is diagonalizable over \(\mathbb F\).
\end{corollary}

The criterion does not require computing the dimensions of the individual
eigenspaces. When the eigenspaces do not span the whole space, there is
no eigenbasis. The nonzero nilpotent matrix in
\autoref{ex:linear-algebra-i:same-polynomial-not-similar} is an example.

\subsection{Approximation by diagonalizable matrices}

\begin{proposition}
  \label{thm:linear-algebra-i:density-of-diagonalizable-complex-matrices}
  The matrices diagonalizable over \(\C\) are dense in
  \(\mathrm M_{n\times n}(\C)\).
\end{proposition}

Here density means that every complex matrix can be approximated
arbitrarily closely, entry by entry, by matrices diagonalizable over
\(\C\).
% October 22 p. 3 describes density in every matrix neighbourhood.
% On October 28 Lucas clarified that this is a statement over C.
% No proof is supplied; do not identify density with a full-measure claim.

